% =====================================================================
%  Chapter 27 — Bistatic Radar and the Australian Record
% =====================================================================

\chapter{Bistatic Radar and the Australian Record}
\label{ch:bistatic-radar}

Radar is the instrument that closes the whole accountability chain, and the
strongest demonstration of it in the world is a pair of Australian radio
telescopes. This chapter takes the technique at its best, because a
comparison that only examines the competitors at their weakest is not a
comparison.

\section{What bistatic radar buys}

A single radar dish transmits and receives. The echo's delay gives the
range to the target and its Doppler shift gives the line-of-sight velocity;
together they reconstruct a three-dimensional shape, and the radar
cross-section gives the size.

A bistatic system separates the two functions. One station transmits and
another receives, and the geometry between them is the baseline. The
advantages are a much longer effective baseline, a different signal-to-noise
regime, and --- importantly for this book --- the possibility that one of the
two stations is not a radio telescope at all.

\section{The Australian record}

The demonstration record is the \idxw{Green Bank Telescope} and the \idxw{Parkes}
observatory used as a bistatic pair, with the Effelsberg 100-metre telescope
contributing in some campaigns. More than fifty near-Earth objects have been
characterised since 2015, including delay-Doppler images and polarimetric
measurements.

The polarimetry is the part that matters most for this book, because it is a
direct diagnostic of the surface state and therefore of the density question
that Chapter 7 identified as the second half of the energy gap. A radar
echo that is polarised in a particular way is a surface that is consolidated
rather than particulate, and that is a statement about the bulk density.

The delay-Doppler images are the part that matters for the shape, because
they are three-dimensional without a convexity assumption.

\begin{table}[!htbp]
\centering
\small
\caption{What the bistatic record delivers, and what limits it.}
\label{tab:bistatic}
\begin{tabularx}{\textwidth}{@{}P{2.57in}Y@{}}
\toprule
\rowcolor{panel}
\textbf{Quantity} & \textbf{Status} \\
\midrule
Objects characterised & more than fifty since 2015 \\
\rowcolor{panel}
Shape & three-dimensional, from delay-Doppler inversion, with no convexity
assumption \\
Polarimetry & a direct surface-state diagnostic and therefore a density
diagnostic \\
\rowcolor{panel}
Selection & only radar-crossing objects, and only at close range \\
Sensitivity growth & three times the Green Bank sensitivity more than
doubles the accessible volume, extending the range by about thirty per
cent; five times more than triples it, by about fifty per cent \\
\bottomrule
\end{tabularx}
\end{table}

The sensitivity line in \tabref{tab:bistatic} is worth pausing on, because
it shows the shape of the problem. Doubling the sensitivity extends the
range by thirty per cent; quintupling it extends the range by fifty per
cent. The relationship is logarithmic, because the fourth-power law means
that the range grows as the fourth root of the sensitivity. **No amount of
telescope improvement moves radar from a population of dozens to a
population of thousands.** A factor of five in sensitivity buys a factor of
1.5 in range, and the range is what selects the objects.

\begin{deeper}{The fourth-power law in one sentence}
The echo power falls as the inverse fourth power of range, so the range
grows as the fourth root of the sensitivity. Improving a radar system by a
factor of a hundred in sensitivity extends its reach by a factor of three.
The population it can reach is therefore set by geometry, not by
engineering, and this is why radar is a complement to the optical
programme and not a replacement for it.
\end{deeper}

\section{The Moon as a bistatic partner}

The lunar bistatic radar concept has been in the literature since 2018, and
the conference literature of 2022 lists ``using the Moon itself'' among the
demonstrated reference-receiver scenarios. The concept is that a radar
transmitter on the Earth and a receiver on the Moon, or the reverse, forms a
bistatic pair with the Moon's surface as the reflector.

The advantage is the same one that motivates the optical case, and it is
worth stating in the same terms: the Moon's surface is a stable, rigid,
airless reflector at a known position, and a bistatic baseline to it is long
and exactly known.

The disadvantage is the one the fourth-power law imposes on everything: the
geometry that makes the baseline long also makes the round trip long, and
the target's radar cross-section has to survive the whole path. The concept
is a good one and the book records it as prior art rather than as a
contribution.

\section{Why radar cannot be the answer to the catalogue}


The conclusion of this chapter is not that radar is unimportant. Radar is
the only method in the field that delivers the mass and the density, and
those are the two terms the accounting chain cannot do without.

The conclusion is that radar cannot be the answer to the catalogue, for
three reasons that are all structural.

The selection effect is the first. Only objects whose surfaces are
radar-bright are observable at all, and for the small near-Earth class that
is a small fraction.

The geometry is the second. The fourth-power law fixes the reachable
population, and the sensitivity scaling is logarithmic in range, so
engineering effort does not move the population.

The opportunity is the third. A radar observation requires the object to be
at a particular distance at a particular time, with a particular geometry,
and the opportunity cannot be scheduled. It has to be caught.

\begin{table}[!htbp]
\centering
\small
\caption{Radar and the optical array against the five jobs the accountability
chain requires.}
\label{tab:radar-vs-optical}
\begin{tabularx}{\textwidth}{@{}P{1.35in}P{1.21in}Y@{}}
\toprule
\rowcolor{panel}
\textbf{Job} & \textbf{Radar} & \textbf{Lunar optical array} \\
\midrule
Diameter & yes, with a model & yes, without a model \\
\rowcolor{panel}
Shape & yes, three-dimensional & yes, projected, with a model \\
Density & yes, by polarimetry & only at a close encounter, by the mass \\
\rowcolor{panel}
Mass & yes & only at a close encounter \\
Opportunity & cannot be scheduled & can be scheduled \\
\rowcolor{panel}
Population reached & dozens & about 37 per year, on request \\
\bottomrule
\end{tabularx}
\end{table}

\tabref{tab:radar-vs-optical} is the comparison that the portfolio decision
of the next chapter turns on. Radar wins five of the six rows. It loses the
last one, and the last one is the row that decides whether a capability can
be commissioned for a catalogue.

\begin{keytakeaways}
\item The Green Bank and Parkes\idx{Parkes} bistatic pair has characterised more than
fifty near-Earth objects since 2015, with three-dimensional shapes and
polarimetry that is a direct density diagnostic.
\item Sensitivity improvements scale as the fourth root in range, so a factor
of five in sensitivity extends the reach by fifty per cent and cannot move
radar from dozens of objects to thousands.
\item A lunar bistatic partner is well documented prior art, and the Moon's
usefulness as a reflector is the same stability argument as the optical
case.
\item Radar wins five of the six jobs the accountability chain requires, and
loses the one that decides whether a capability can serve a catalogue.
\end{keytakeaways}
